\begin{resumo}
Este trabalho apresenta a proposta e o desenvolvimento inicial de um Avaliador de
Serviços Públicos voltado a serviços digitais, no contexto da transformação
digital do Estado brasileiro e da necessidade de tornar a avaliação da
experiência do cidadão mais útil para a melhoria contínua. O problema investigado
parte da baixa utilidade de avaliações exibidas apenas ao final da jornada, pois
muitos usuários encontram dificuldade antes de concluir o serviço e nem sempre
respondem a formulários tradicionais. O objetivo geral é propor e prototipar uma
solução capaz de registrar sinais de dificuldade na navegação, sugerir
redirecionamento em serviços prioritários, solicitar feedback contextual de baixo
esforço e apresentar indicadores agregados para apoio à melhoria de serviços
públicos digitais. A metodologia combinou revisão bibliográfica e normativa,
Design Science Research, prototipação incremental e Lean Inception. A partir das
dinâmicas de visão de produto, delimitação de escopo, objetivos, personas,
jornadas, brainstorming de funcionalidades, revisão técnica, sequenciamento e
Canvas MVP, foi definido um MVP composto por sessão técnica anônima, registro de
cliques e tempo, detecção por regras simples, feedback contextual, formulário
estático comparativo e painel administrativo. O protótipo foi implementado com
uma API local, páginas simuladas de portal público, formulário de avaliação por
estrelas, atributos, NPS e comentário, além de dados armazenados em SQLite. Como
resultado, o trabalho demonstra a viabilidade conceitual e técnica de uma solução
que transforma eventos de navegação e feedback declarado em evidências
administrativas. Conclui-se que o feedback contextual pode complementar
instrumentos oficiais de avaliação, desde que usado com finalidade de melhoria,
cuidado com privacidade, governança institucional dos dados e validação futura
com usuários reais.

\vspace{\onelineskip}

\noindent
\textbf{Palavras-chave}: serviços públicos. feedback contextual. governo
digital. experiência do usuário. protótipo.
\end{resumo}
